Dans une station balnéaire on a interrogé 600 touristes, français ou étrangers, sur leur séjour.

Tous ont répondu. Les effectifs sont reportés dans le tableau suivant :

\begin{center}
\begin{tblr}{
	colspec = {l|c|c|c|c},
	hlines, vlines,
	row{1} = {font=\bfseries, bg=gray!15},
	column{1} = {font=\bfseries, bg=gray!15}
}
 & Camping & Location & Hôtel & Total \\
Français & 232 & 116 & 12 & 360 \\
Étrangers & 96 & 96 & 48 & 240 \\
Total & 328 & 212 & 60 & 600 \\
\end{tblr}
\end{center}

Les réponses des questions suivantes seront données sous forme de fractions non nécessairement simplifiées.

On choisit une personne au hasard. On suppose que toutes les personnes ont la même probabilité d'être choisies.

On considère les événements :
\begin{itemize}[left=5mm]
	\item $E$ : « La personne choisie est un touriste étranger » ;
	\item $C$ : « La personne choisie séjourne dans un camping ».
\end{itemize}

\begin{enumerate}
	\item Donner la probabilité $P(C)$ de l'événement $C$.
	\item Donner la probabilité $P(\overline{C})$ où $\overline{C}$ désigne l'événement contraire de $C$.
	\item Donner la probabilité de l'événement : « La personne interrogée est un touriste étranger et séjourne dans un camping ».
	\item Sachant que la personne choisie est en location, donner la probabilité qu'elle soit un touriste français.
\end{enumerate}
